% =============================================================
% template.tex  —  LaTeX
%
% FLAGS (set these before \begin{document}):
%
%   \binding   : left | right   (which physical side is the spine)
%   \sides     : single | double (one-sided or alternating pages)
%   \notesside : inner | outer   (which side margin notes appear on)
%
% MARGIN WIDTH (set \marginwidth):
%   absolute:  \setlength{\marginwidth}{3cm}
%   relative:  \setlength{\marginwidth}{0.2\paperwidth}
%
% HEAD / FOOT HOOKS (set these in your document):
%   \headinner, \headmiddle, \headouter
%   \footinner, \footmiddle, \footouter
%
%   In double-sided mode these alternate physically with each page.
%   In single-sided mode inner=left, outer=right throughout.
%
% MARGIN NOTES:
%   \marginnote{REF}{NOTE TEXT}
%   REF is any user-chosen symbol (★, †, 42, …).
%   Deposits REF inline and REF + NOTE in the margin.
%   Placement is approximate; REF is the binding reference.
% =============================================================

% --- sides flag drives the documentclass option --------------
% Set \sides before loading the class by using a wrapper flag:
%   double-sided:  \documentclass[a4paper,11pt,twoside]{article}
%   single-sided:  \documentclass[a4paper,11pt,oneside]{article}
%
% We use twoside here; change to oneside for single-sided output.
\documentclass[a4paper, 11pt, twoside]{article}

\usepackage{geometry}
\usepackage{marginfix}
\usepackage{fancyhdr}
\usepackage{ragged2e}
\usepackage{ifthen}

% =============================================================
% FLAGS
% =============================================================

\newcommand{\binding}{left}    % left | right
\newcommand{\sides}{double}    % single | double
\newcommand{\notesside}{outer} % inner | outer

% =============================================================
% MARGIN WIDTH
% =============================================================

\newlength{\marginwidth}
\setlength{\marginwidth}{0.2\paperwidth}  % RELATIVE — change to e.g. 3cm for absolute

% =============================================================
% GEOMETRY
% Resolves \binding, \sides, and \notesside into geometry params.
%
% LaTeX's geometry package uses "inner" (binding side) and "outer"
% (fore-edge side) natively in twoside mode, which matches our
% logical naming exactly. In oneside mode inner=left, outer=right.
%
% marginpar is placed on the \notesside. We achieve this by
% setting marginparwidth on the appropriate side and using
% \reversemarginpar when \notesside=inner.
% =============================================================

\ifthenelse{\equal{\notesside}{outer}}{%
  \geometry{
    a4paper,
    top=2.5cm, bottom=2.5cm,
    inner=2.5cm, outer=2.5cm,
    marginparwidth=\marginwidth,
    marginparsep=0.5em,
  }%
}{%
  % notesside=inner: swap marginpar to the binding side
  \geometry{
    a4paper,
    top=2.5cm, bottom=2.5cm,
    inner=2.5cm, outer=2.5cm,
    marginparwidth=\marginwidth,
    marginparsep=0.5em,
  }%
  \reversemarginpar  % places \marginpar on the inner (binding) side
}

% binding=right means the spine is on the right, so we swap
% inner and outer by reversing the document direction:
\ifthenelse{\equal{\binding}{right}}{%
  \reversemarginpar  % combined with geometry's twoside, this flips binding side
  % Note: for full RTL support, consider the 'mirror' option in geometry.
}{}

% =============================================================
% HEAD / FOOT HOOKS
% =============================================================

% User-facing logical hooks — set these in your document:
\newcommand{\headinner}{}
\newcommand{\headmiddle}{}
\newcommand{\headouter}{\thepage}

\newcommand{\footinner}{}
\newcommand{\footmiddle}{}
\newcommand{\footouter}{}

% Wire logical hooks to fancyhdr physical slots.
%
% In twoside mode fancyhdr uses:
%   [LE] = left on even (verso) pages  = inner on verso = binding side
%   [RO] = right on odd (recto) pages  = inner on recto = binding side
%   [RE] = right on even pages         = outer on verso
%   [LO] = left on odd pages           = outer on recto
%
% So inner  -> [LE,RO]  and  outer -> [RE,LO]
% In oneside mode we use [L] and [R] directly.

\pagestyle{fancy}
\fancyhf{}

\ifthenelse{\equal{\sides}{double}}{%
  \fancyhead[LE,RO]{\headinner}
  \fancyhead[C]{\headmiddle}
  \fancyhead[RE,LO]{\headouter}
  \fancyfoot[LE,RO]{\footinner}
  \fancyfoot[C]{\footmiddle}
  \fancyfoot[RE,LO]{\footouter}
}{%
  % single-sided: inner=left, outer=right
  \fancyhead[L]{\headinner}
  \fancyhead[C]{\headmiddle}
  \fancyhead[R]{\headouter}
  \fancyfoot[L]{\footinner}
  \fancyfoot[C]{\footmiddle}
  \fancyfoot[R]{\footouter}
}

\renewcommand{\headrulewidth}{0.4pt}
\renewcommand{\footrulewidth}{0pt}

% =============================================================
% MARGIN NOTE MACRO
%
% \marginnote{REF}{NOTE}
%   REF  — user-supplied reference symbol, placed inline and in margin
%   NOTE — note text, set smaller in the margin column
%
% Uses \marginpar for placement (float-like, but REF binds the pair).
% \reversemarginpar above already directs \marginpar to the correct side.
% =============================================================

\newcommand{\marginnote}[2]{%
  \mbox{}\textsuperscript{#1}%                   % inline reference mark
  \marginpar{%
    \footnotesize\RaggedRight
    \textsuperscript{#1}\,#2%                     % margin: mark + note text
  }%
}

% =============================================================
% DOCUMENT
% =============================================================
\begin{document}

% --- Set flags (can be changed here before any content) ------
\renewcommand{\binding}{left}
\renewcommand{\sides}{double}
\renewcommand{\notesside}{outer}

% --- Set head/foot content -----------------------------------
\renewcommand{\headinner}{\thepage}
\renewcommand{\headmiddle}{}
\renewcommand{\headouter}{My Document}

% --- Body text -----------------------------------------------
Lorem ipsum dolor sit amet, consectetur adipiscing elit.%
\marginnote{★}{This note is keyed to the star in the body text.}
Sed do eiusmod tempor incididunt ut labore et dolore magna aliqua.

Ut enim ad minim veniam, quis nostrud exercitation ullamco laboris.%
\marginnote{†}{A dagger marks this second note.}
Duis aute irure dolor in reprehenderit in voluptate velit esse cillum.

\end{document}
